\section{Introduction}
\label{sec:intro}

Two questions follow a change to a tax or a benefit, and they are not the same question. The first is arithmetic: given this family, these earnings, these children, this year, what does the rule book say they pay and receive, before and after? The second is statistical: across a whole country, how much does the change raise, and who gains and loses? The first has an exact answer, because tax and benefit law is a computable function of a household's circumstances and the answer can be checked with a pencil. The second does not, because nobody observes the whole country: it is estimated from a survey that is sampled, imputed, reweighted and aged forward, and every one of those steps puts error into the headline number.

This paper documents the PolicyEngine tax--benefit microsimulation---the engine behind \url{https://policyengine.org}, and the lead member of the \emph{PolicyEngine Macro} suite---as this suite uses it, and it treats the distinction between those two kinds of number as its subject.\footnote{\url{https://policyengine-macro.vercel.app/pe/}. PolicyEngine Macro is a suite of open economic models---microsimulation, macroeconometric, structural-VAR, heterogeneous-agent, overlapping-generations and ecological stock-flow---that score the same PolicyEngine reform object and report comparable real-world aggregates \citep{policyengine}.} The microsimulation is the suite's first model and its only production-ready member, in the specific sense the capability registry records: ``production-ready for selected household applications.'' It is also the only member covering both the United Kingdom and the United States, and the only one whose evidence base is implemented legislation rather than another institution's published output. Every bridge in the suite starts or ends here: it supplies the static costing that the OBR macroeconometric emulator turns into second-round demand effects \citep{ahmadi2026obr}, the tax functions that the overlapping-generations member estimates its household problem from \citep{ahmadi2026og}, and the incidence engine through which macro paths from the other members are pushed back down to households.

The contribution is threefold. First, \emph{a documented exactness boundary}. The suite's capability registry records this model's uncertainty as ``none for household arithmetic; survey/calibration uncertainty for population estimates,'' and Section~\ref{sec:uncertainty} sets out precisely why the two clauses are different in kind rather than in degree: the arithmetic inside a population aggregate is exactly as exact as a household calculation, and all of the uncertainty lives in the weights and in the imputed components of the inputs. Blurring the two---quoting a population costing with the confidence one is entitled to have in a household calculation---is the characteristic misuse of a microsimulation model, and the boundary is worth stating in a document rather than in a docstring.

Second, \emph{an honest validation posture for a model that does not forecast}. There is no forecast error to report here, because nothing is being forecast: published rules are applied, and the check is whether they were transcribed correctly. Section~\ref{sec:validation} states what that check does and does not establish, and reports the single committed external benchmark---the static costing of 1p on the UK basic rate against HMRC's published ready reckoner---with its actual numbers and its actual gap, including the fact that it is one reform, in one country, on one tax, and that the official figure it is measured against is itself an estimate on a different data basis.

Third, \emph{a declared account of what the model cannot see and what the suite substitutes}. The registry's \texttt{cannot\_answer} list is short and absolute: GDP, inflation, interest rates, macro feedback. Section~\ref{sec:bridge} describes the machinery that supplies those, and is equally explicit about that machinery's limits---in particular that the route into the OBR emulator injects an entire annual costing as a held add-factor on nominal household disposable income, flat within the year, and is therefore a demand-side approximation rather than a general model of reform incidence. The macro model never learns which households paid.

Two boundaries on the paper itself should be stated at the outset. This is a description of the model \emph{as the suite consumes it}: the country rule packages (\texttt{policyengine-uk}, \texttt{policyengine-us}) and the simulation framework beneath them (\texttt{policyengine-core}, forked from OpenFisca) are maintained upstream, and nothing here audits how completely they cover UK or US statute or how thorough their own test suites are. And every figure quoted below comes from an artifact committed to the PolicyEngine Macro repository---the capability registry, the integration adapters, the committed chart data behind the site's validation figures, or the sister working papers---with the source named at the point of use; no number in this paper was produced for it.

The rest of the paper proceeds as follows. Section~\ref{sec:literature} places the model in the microsimulation tradition and in the open-model literature. Section~\ref{sec:model} sets out what the model computes: variables and dated parameters as a dependency graph, the country-specific entity structures, and the reform object. Section~\ref{sec:method} describes the integration layer this suite drives it through. Section~\ref{sec:uncertainty} is the paper's central section, on the boundary between exact arithmetic and estimated aggregates. Section~\ref{sec:validation} gives the validation posture and the ready-reckoner benchmark; Section~\ref{sec:bridge} the limits and the bridges; Section~\ref{sec:repro} reproducibility, data vintage and the gated microdata; Section~\ref{sec:limitations} limitations, and Section~\ref{sec:conclusion} concludes. Appendices give the curated reform-parameter catalogue, the worked household arithmetic, and the reform-object grammar with its refusals.
